\documentclass[10pt,a4paper]{amsart}
\usepackage[margin=19mm]{geometry}
\usepackage{amsmath,amssymb,mathtools}
\usepackage[scr=rsfs,cal=euler]{mathalfa}
\usepackage{xfrac}
\usepackage[unicode,hidelinks,bookmarks=false]{hyperref}
\newcommand{\ie}{i.e.\ }
\newcommand{\cf}{cf.\ }
\newcommand{\resp}{respectively\ }
\newcommand{\Subsection}{\subsection}
\providecommand{\nobreakdash}{\nobreak\hbox{-}}
\setlength{\emergencystretch}{2em}
\multlinegap0pt
\usepackage{amssymb}
\usepackage[normalem]{ulem}
\newtheorem{lemma}{Lemma}
\title{Convergence of semilinear parabolic flows with general initial data}
\author{Daniel Restrepo}
\date{Source: arXiv:2603.00748v1 (CC BY 4.0)}
\begin{document}
\maketitle

\noindent\emph{Source and licence.} Convergence of semilinear parabolic flows with general initial data. Version arXiv:2603.00748v1 (CC BY 4.0), distributed by arXiv under the Creative Commons Attribution 4.0 International licence, http://creativecommons.org/licenses/by/4.0/, as recorded in the arXiv licence metadata for this exact version and shown on the abstract page https://arxiv.org/abs/2603.00748v1. Full licence notice: \texttt{licenses/arxiv-2603.00748-CC-BY-4.0.txt}.\par\smallskip

\noindent\textit{Editorial scope.} Counted unit: the article's lemma ch (source0.tex lines 1619--1643). The article's complete original proof of it is delivered with it (source0.tex lines 1645--1787, one \texttt{proof} environment of 143 lines). No mathematics is written, repaired or re-derived: the statement and the proof are the article's own lines, in the article's own notation, and no retained line is substituted or re-flowed.  No further source-local context is needed: every cross-reference of every retained line resolves inside the two delivered spans.  Nothing was omitted from any retained span, so the package carries no omission record.  No retained line contains a citation, so the wrapper carries no bibliography and no bibliography entry.  No retained line contains a figure environment, an image-inclusion command or drawing code, so no image is delivered and the package declares no asset dependency.  Editorial formatting only, in addition: an AMS \texttt{amsart} preamble that restates the article's own declarations and macros used by the retained lines, namely the environments \texttt{lemma} on separate counters, together with the standard packages those lines need.  The retained environments print in this document as Lemma 1 in order of appearance, whereas the article prints the same items with the numbering of its own full document.  The complete native source of the article as distributed in the arXiv e-print for this exact version is bundled unchanged as \texttt{sources/195505/source0.tex}.
\begin{lemma}\label{lemma ch}
Let \(P=\{x_1,\dots,x_M\}\subset \mathbb{R}^n\) be a collection of distinct points. 
Then there exists a constant \(D=D(M,n)>0\) and a point \(y\in P\) for which one can find a unit vector \(e\in \mathbb{S}^{n-1}\) satisfying
\begin{equation}\label{eq hyper bound}
(x_i-y)\cdot e 
\ge 
\frac{1}{D(M,n)}\,|x_i-y|,
\qquad 
\text{for every } x_i\in P\setminus\{y\}.
\end{equation}
Moreover, if the points of \(P\) satisfy the separation condition
\[
|x_i-x_j|\ge L 
\qquad \text{for all } i\ne j,
\]
then there exists a constant \(D_2=D_2(M,n)>0\) such that
\begin{equation}\label{eq hyper bound neigh}
(x_i-z)\cdot e 
\ge 
\frac{1}{D_2(M,n)}\,|x_i-z|,
\qquad 
\text{for every } x_i\in P\setminus\{y\},
\end{equation}
whenever \(z\in B_{L'}(y)\), where \(L'=\frac{L}{2D(M,n)}\).
\end{lemma}

\begin{proof}
We first prove \eqref{eq hyper bound} by induction on \(M\). If \(M=2\), the statement is immediate: taking
\[
e=\frac{x_2-x_1}{|x_2-x_1|},
\qquad
y=x_1,
\]
the inequality holds with \(D=1\).

Assume now that the statement holds for any collection of \(M-1\) points.  
Apply the inductive hypothesis to \(P\setminus\{x_M\}\). Then there exist 
\(y_0\in P\setminus\{x_M\}\) and \(e_0\in\mathbb S^{n-1}\) such that
\begin{equation}\label{eq hyper bound0}
(x-y_0)\cdot e_0
\ge
\frac{1}{D(M-1,n)}|x-y_0|,
\qquad
\text{for all } x\in P\setminus\{x_M,y_0\}.
\end{equation}

Set
\[
K=\max\{4D(M-1,n),2\}.
\]
We distinguish two cases.

\medskip
\noindent
\textit{Case 1:} Assume
\begin{equation}\label{eq noortho}
\frac{|(x_M-y_0)\cdot e_0|}{|x_M-y_0|}
\ge
\frac{1}{K}.
\end{equation}

If \((x_M-y_0)\cdot e_0>0\), then \eqref{eq hyper bound} holds with 
\(y=y_0\), \(e=e_0\), and \(D(M,n)=K\).

If \((x_M-y_0)\cdot e_0<0\), we set \(y=x_M\) and \(e=e_0\).  
From \eqref{eq noortho} we obtain
\[
(y_0-x_M)\cdot e_0
\ge
\frac{1}{K}|x_M-y_0|.
\]
For \(x\in P\setminus\{x_M,y_0\}\), we deduce from \eqref{eq hyper bound0} that
\[
(x-x_M)\cdot e_0
=
(x-y_0)\cdot e_0+(y_0-x_M)\cdot e_0
\ge
\frac{1}{D(M-1,n)}|x-y_0|
+
\frac{1}{K}|x_M-y_0|.
\]
Using the triangle inequality and the definition of \(K\), we deduce
\[
(x-x_M)\cdot e_0
\ge
\frac{1}{K}|x-x_M|.
\]

\medskip
\noindent
\textit{Case 2:} Assume
\begin{equation}\label{eq aortho}
\frac{|(x_M-y_0)\cdot e_0|}{|x_M-y_0|}
<
\frac{1}{K}.
\end{equation}
Let
\[
e_1=\frac{x_M-y_0}{|x_M-y_0|},
\qquad
\alpha=\frac{2}{K},
\qquad
e=\frac{e_0+\alpha e_1}{|e_0+\alpha e_1|}.
\]
Then \(|e_0\cdot e_1|\le 1/K\), and therefore
\begin{equation}\label{eq e1e2}
|e_0+\alpha e_1|^2
\le
1+\frac{8}{K^2}
\le 4.
\end{equation}

A direct computation gives
\[
(x_M-y_0)\cdot e
=
|x_M-y_0|
\frac{e_0\cdot e_1+\alpha}{|e_0+\alpha e_1|}
\ge
\frac{1}{2K}|x_M-y_0|.
\]

If \(x\in P\setminus\{x_M,y_0\}\), writing 
\(\nu_x=\frac{x-y_0}{|x-y_0|}\), the inductive hypothesis gives
\(\nu_x\cdot e_0\ge 1/D(M-1,n)\). Hence
\[
(x-y_0)\cdot e
=
|x-y_0|
\frac{e_0\cdot\nu_x+\alpha e_1\cdot\nu_x}{|e_0+\alpha e_1|}
\ge |x-y_0|
\frac{1/D(M-1,n)-\alpha }{2} \ge
\frac{1}{4D(M-1,n)}|x-y_0|.
\]

This completes the proof of \eqref{eq hyper bound}.

\medskip

We now prove \eqref{eq hyper bound neigh}.  
Let \(z\in B_{L'}(y)\) and write \(z=y+z_0\) with \(|z_0|\le L'\).  
For \(x\in P\setminus\{y\}\),
\[
(x-z)\cdot e
=
(x-y)\cdot e - z_0\cdot e
\ge
\frac{1}{D(M,n)}|x-y|-|z_0|.
\]
Since \(|x-y|\ge L\) and 
\(L'=\frac{L}{2D(M,n)}\), we obtain
\[
(x-z)\cdot e
\ge
\frac{1}{2D(M,n)}|x-y|.
\]
Finally,
\[
|x-z|
\le
|x-y|+|z-y|
\le
|x-y|+L'
\le
\Bigl(1+\frac{1}{2D(M,n)}\Bigr)|x-y|,
\]
which yields \eqref{eq hyper bound neigh} with a suitable constant 
\(D_2(M,n)>0\).
\end{proof}

\end{document}
